\documentclass[11pt,a4paper]{article}

\usepackage[utf8]{inputenc}
\usepackage[T1]{fontenc}
\usepackage[spanish]{babel}
\usepackage[margin=2.5cm]{geometry}
\usepackage{lmodern}
\usepackage[table]{xcolor}
\usepackage{array}
\usepackage{tabularx}
\usepackage{enumitem}
\usepackage{fancyhdr}
\usepackage[hidelinks]{hyperref}

\definecolor{azul}{HTML}{17365D}
\definecolor{azulclaro}{HTML}{EDF2F7}

\setlength{\parindent}{0pt}
\setlength{\parskip}{6pt}
\setlength{\headheight}{14pt}
\setlength{\tabcolsep}{10pt}
\renewcommand{\arraystretch}{1.5}

\setlist[enumerate]{
    leftmargin=*,
    itemsep=0.6em
}

\pagestyle{fancy}
\fancyhf{}
\fancyhead[L]{\small UTEQ · Ingeniería de Requerimientos}
\fancyhead[R]{\small GA-TL-02}
\fancyfoot[L]{\footnotesize Grupo CEV}
\fancyfoot[R]{\footnotesize Página \thepage}
\renewcommand{\headrulewidth}{0.4pt}
\renewcommand{\footrulewidth}{0.4pt}

\begin{document}

\begin{center}

\vspace*{0.3cm}

{\large\bfseries
UNIVERSIDAD TÉCNICA ESTATAL DE QUEVEDO}

\vspace{0.5cm}

{\small
Facultad de Ciencias de la Computación · Carrera de Software}

\vspace{0.5cm}

{\bfseries ACTIVIDAD FORMATIVA GRUPAL}

\vspace{0.4cm}

{\large\bfseries GA-TL-02 -- FORM}

\vspace{0.4cm}

{\large\bfseries
Taller de entrevista estructurada y design thinking
aplicados al PFC}

\end{center}

\vspace{0.4cm}

\section*{Datos informativos}

\arrayrulecolor{azul}

\begin{tabularx}{\textwidth}{
    |>{\columncolor{azulclaro}\raggedright\arraybackslash}p{4.1cm}
    |>{\raggedright\arraybackslash}X|
}
\hline
\rowcolor{azul}
\textcolor{white}{\textbf{Campo}} &
\textcolor{white}{\textbf{Detalle}} \\
\hline
Integrantes · CEV &
Contreras Chavez Kevin German \newline
Espinoza Pivaque Angelo Maximiliano \newline
Vaca Romero David Octavio \\
\hline
Asignatura &
Ingeniería de Requerimientos (ISR-401) \\
\hline
Nivel y paralelo &
Cuarto nivel · Paralelo A \\
\hline
Docente &
Gleiston Cicerón Guerrero Ulloa, PhD \\
\hline
Unidad &
Unidad 2. Elicitación y Análisis de Requisitos \\
\hline
Proyecto &
Sistema para la identificación temprana de dificultades
académicas mediante el análisis de indicadores de
rendimiento estudiantil \\
\hline
Fecha límite &
Sábado 21 de noviembre de 2026, 23:59 \\
\hline
Período &
2026--2027 SPA \\
\hline
\end{tabularx}

\vspace{0.4cm}

\textbf{Repositorio del equipo}

{\small
\mbox{\url{https://github.com/Espinoza-A-1006/ISR401-CEV}}
}

\newpage

\section{Contexto y propósito del ensayo}

Este ensayo tiene como propósito practicar la aplicación de
una entrevista estructurada mediante un juego de roles entre
compañeros, tomando como referencia el Sistema para la
identificación temprana de dificultades académicas mediante
el análisis de indicadores de rendimiento estudiantil.
La entrevista se concentra en las tutorías como una forma
de acompañamiento ante las dificultades académicas.
Se busca explorar las necesidades del estudiante representado
y evaluar la claridad del guion. La información obtenida
servirá como base para elaborar el mapa de empatía, formular
el problema y proponer ajustes a las preguntas, dentro del
alcance de una simulación académica.

\subsection{Contexto del proyecto}

El proyecto consiste en un sistema orientado a la identificación
temprana de dificultades académicas mediante el análisis de
indicadores de rendimiento estudiantil \cite{chang2025, martins2019}. Las fuentes de información,
los indicadores y los criterios de interpretación deberán
definirse con los responsables académicos. En este ensayo se
aborda específicamente la experiencia del estudiante al solicitar
y recibir tutorías como una forma de apoyo académico.

Como punto de partida, el equipo plantea un proceso en el que los
estudiantes contactan presencialmente al docente para comunicar sus
dudas y solicitar apoyo. Cuando se acuerda realizar una tutoría,
la coordinación continúa mediante WhatsApp, donde se definen el día
y la hora, se confirma la sesión y se comunican posibles cambios.

La información distribuida entre conversaciones presenciales y
mensajes puede dificultar la consulta de los acuerdos establecidos.
Cuando se proponen varios horarios, los participantes necesitan
identificar cuál fue confirmado y qué solicitudes continúan
pendientes.

El seguimiento posterior también requiere recuperar los temas
revisados, las recomendaciones recibidas y las actividades pendientes.
Disponer de estos antecedentes permite retomar el acompañamiento
en sesiones posteriores.

Este contexto constituye un planteamiento inicial para orientar
el ensayo. Las respuestas obtenidas permitirán explorar las
situaciones descritas y revisar la utilidad del guion, considerando
que proceden de una simulación académica. Los resultados se limitarán
a la experiencia tutorial explorada y no permitirán, por sí solos,
definir o validar los indicadores de rendimiento del sistema.

\subsection{Objetivo de la entrevista}

Explorar cómo un estudiante solicita, coordina y da continuidad
a sus tutorías académicas como forma de apoyo ante sus dificultades,
para identificar necesidades de información antes y después de
las sesiones. Asimismo, evaluar la claridad de las preguntas y
su utilidad para obtener información pertinente.

\subsection{Modalidad y rol del participante}

La actividad se desarrollará mediante un juego de roles entre
compañeros. Kevin German Contreras Chavez asumirá el rol de
entrevistador y Angelo Maximiliano Espinoza Pivaque representará
a un estudiante que solicita tutorías académicas.

El ensayo permitirá practicar la entrevista estructurada e
identificar posibles ajustes en el guion. Las respuestas y
los hallazgos se presentarán como resultados de una simulación
académica.

\section{Guion de entrevista estructurada y respuestas de apoyo (P01)}

El guion consta de una apertura, trece preguntas organizadas
en cuatro bloques temáticos y un cierre. Las preguntas abordan
la solicitud, coordinación, confirmación y seguimiento de las
tutorías académicas, dentro del contexto de apoyo al estudiante
ante sus dificultades de aprendizaje. Se registran junto con las
respuestas de apoyo preparadas para el juego de roles del
participante \texttt{P01}.

\subsection{Apertura}

Hola, gracias por participar. Realizaré esta entrevista como parte del taller GA-TL-02
de Ingeniería de Requerimientos. En este ensayo representarás
a un estudiante que solicita tutorías académicas.

Nuestro proyecto se orienta a la identificación temprana de
dificultades académicas mediante el análisis de indicadores
de rendimiento estudiantil. En esta entrevista nos centraremos
en las tutorías como una forma de apoyo: queremos conocer
cómo solicitas y coordinas las sesiones, qué dificultades
encuentras y cómo recuperas lo trabajado para continuar
tu aprendizaje. También revisaremos si las preguntas
del guion son claras y permiten obtener información útil.
La entrevista contiene 13 preguntas y tendrá una duración
aproximada de 10 minutos, dependiendo de la extensión de tus
respuestas.

Tu participación es voluntaria. No hay respuestas correctas
o incorrectas y no necesitas inventar información: puedes indicar
si una situación no te ha ocurrido o si no sabes cómo responder.
Puedes solicitar aclaraciones, omitir cualquier pregunta, pedir
una pausa o finalizar tu participación sin necesidad de justificarlo.

Con tu autorización, grabaremos únicamente el audio para revisar
las respuestas y analizar el desarrollo de la entrevista.
La grabación se conservará en un espacio privado con acceso
restringido al equipo y al docente. No se publicará en GitHub.
En el informe público utilizaremos el código P01 para identificar
tus respuestas y evitaremos incluir datos personales innecesarios.
Los resultados se presentarán expresamente como información
obtenida en una simulación académica.

Antes de comenzar, ¿tienes alguna pregunta sobre la actividad
o el uso de la información?

¿Aceptas participar voluntariamente en este ensayo?

¿Autorizas la grabación del audio para los fines académicos
que te he explicado?

\subsection{Bloque 1: Solicitud y coordinación}

\begin{enumerate}[label=\textbf{\arabic*.}]
    \item ¿Cómo solicitas actualmente una tutoría cuando tienes
    dudas sobre alguna materia?

    \textit{Respuesta de apoyo (P01):} Normalmente espero a que termine la clase y me acerco al docente. Le explico qué parte no entendí y le pregunto cuándo podría ayudarme a revisarla.

    \item ¿Por qué medios te comunicas con el docente para
    iniciar el contacto de una tutoría?

    \textit{Respuesta de apoyo (P01):} Primero hablo con él en persona. Después, si quedamos en hacer la tutoría, le escribo por WhatsApp para coordinar el horario.

    \item ¿Cómo se define o acuerda el día y la hora en que
    se realizará la tutoría?

    \textit{Respuesta de apoyo (P01):} El docente me dice cuándo tiene tiempo y reviso si puedo asistir. Si tengo otra clase a esa hora, buscamos otro momento que nos sirva a los dos.

    \item ¿Cómo se gestiona un cambio en el horario acordado
    para una tutoría?

    \textit{Respuesta de apoyo (P01):} Le escribo para avisarle que no puedo asistir y le pregunto si podemos cambiarla. Espero que me confirme el nuevo horario antes de darlo por hecho.
\end{enumerate}

\subsection{Bloque 2: Confirmación e información previa}

\begin{enumerate}[resume,label=\textbf{\arabic*.}]
    \item ¿Cómo identificas actualmente si una tutoría ya
    quedó confirmada en fecha y hora?

    \textit{Respuesta de apoyo (P01):} Cuando el docente me responde que sí puede atenderme ese día y a esa hora. Si solo me dice que quizá tenga tiempo, entiendo que todavía no está confirmada.

    \item Si se proponen varios horarios posibles en una
    conversación, ¿cómo verificas cuál fue el que finalmente
    se acordó?

    \textit{Respuesta de apoyo (P01):} Reviso los últimos mensajes para encontrar el horario que aceptamos. Si no me queda claro, prefiero preguntarle otra vez antes de ir a una hora equivocada.

    \item ¿Qué información necesitas tener clara antes de
    asistir a una sesión de tutoría?

    \textit{Respuesta de apoyo (P01):} Principalmente el día, la hora y dónde nos vamos a reunir. También necesito saber si debo llevar ejercicios o preparar algo para aprovechar mejor el tiempo.
\end{enumerate}

\subsection{Bloque 3: Desarrollo y seguimiento}

\begin{enumerate}[resume,label=\textbf{\arabic*.}]
    \item ¿Cómo conservas la información que revisas durante
    una tutoría?

    \textit{Respuesta de apoyo (P01):} Anoto en mi cuaderno lo que voy entendiendo y guardo los ejercicios que hacemos. Si el docente me envía un archivo, también lo guardo para practicar después.

    \item Después de que termina la tutoría, ¿dónde puedes
    consultar el material utilizado, si lo hubo?

    \textit{Respuesta de apoyo (P01):} En mis apuntes o en los archivos que me haya compartido el docente. A veces tengo que buscar en WhatsApp porque ahí quedó el documento o la explicación.

    \item Si semanas después vuelves a necesitar apoyo en
    la misma materia, ¿cómo recuperas lo trabajado en
    sesiones anteriores?

    \textit{Respuesta de apoyo (P01):} Busco los ejercicios anteriores y reviso mis apuntes. Si no anoté algo o no encuentro el archivo, tengo que volver a preguntarle al docente para recordar dónde nos quedamos.
\end{enumerate}

\subsection{Bloque 4: Dificultades y cierre}

\begin{enumerate}[resume,label=\textbf{\arabic*.}]
    \item De todo el proceso, desde solicitar la tutoría
    hasta repasar lo visto, ¿cuál consideras que es la
    principal dificultad?

    \textit{Respuesta de apoyo (P01):} Para mí sería encontrar después toda la información. Puedo tener una explicación en el cuaderno y otra en WhatsApp, y termino buscando en varios lugares para recordar qué debía practicar.

    \item ¿Cómo influye la duración de la tutoría en tu
    aprovechamiento de la sesión?

    \textit{Respuesta de apoyo (P01):} Depende de qué tan complicado sea el tema. No me basta con escuchar la explicación; también necesito intentar un ejercicio para saber si entendí. Si no alcanza el tiempo para eso, puedo salir todavía con dudas.

    \item ¿Hay algún otro aspecto sobre el acompañamiento
    académico que te gustaría agregar?

    \textit{Respuesta de apoyo (P01):} Me serviría terminar la tutoría teniendo claro qué debo practicar. Y que, cuando vuelva, podamos revisar esos ejercicios para ver si ya superé la dificultad o todavía necesito ayuda.
\end{enumerate}

\subsection{Cierre}

Gracias por compartir tus respuestas. Antes de finalizar,
voy a resumir los aspectos principales que mencionaste
para comprobar que los interpreté correctamente.

\textit{Resumen del entrevistador:} Por lo que me comentaste, primero contactas al docente en persona y después coordinan la tutoría por WhatsApp. Para asistir, necesitas tener claros el horario, el lugar y los materiales que debes preparar. También mencionaste que guardas lo trabajado en apuntes y archivos, pero a veces resulta difícil recuperarlo porque está repartido en distintos lugares. Te ayudaría tener claras las actividades que debes practicar y retomarlas en la siguiente sesión. Además, consideras importante resolver ejercicios para comprobar que comprendiste la explicación.

¿Entendí bien o hay algo que quieras corregir o agregar?

\textit{Respuesta de apoyo (P01):} Sí, eso resume lo que quise decir. Lo que más me ayudaría es tener claro el horario y poder recuperar lo que trabajamos en las tutorías.

Para ayudarnos a mejorar el guion, ¿hubo alguna pregunta
que te resultara confusa, repetitiva o difícil de responder?

\textit{Respuesta de apoyo (P01):} No, la verdad entendí todas las preguntas y pude responderlas.

La información se utilizará para elaborar el mapa de empatía,
formular el problema y mejorar el guion, dentro de esta
simulación académica. El audio permanecerá en el espacio
privado y con el acceso restringido que explicamos al inicio.

Antes de terminar, ¿tienes alguna pregunta sobre el uso
de tus respuestas o de la grabación?

Gracias por tu tiempo y colaboración. Con esto finalizamos
el ensayo y detenemos la grabación.

\section{Registro del ensayo y consentimiento}

\subsection{Fecha, modalidad y participantes}

El ensayo se realizó el 20 de noviembre de 2026,
a las 20:00, mediante una llamada virtual en Zoom.
La grabación de audio tiene una duración de
9 minutos y 24 segundos.

Kevin German Contreras Chavez participó como
entrevistador y Angelo Maximiliano Espinoza Pivaque
como entrevistado, representando a un estudiante
que solicita tutorías académicas.

La actividad se desarrolló como una simulación
académica para aplicar el guion de entrevista
estructurada y explorar las necesidades del
estudiante representado.

\subsection{Consentimiento informado y resguardo de datos}

Al inicio de la sesión se informó el propósito académico de la simulación, el carácter voluntario de la participación y el registro exclusivo en audio bajo el código \texttt{P01}. El participante aceptó voluntariamente y autorizó la grabación del audio para fines académicos y de análisis del equipo y del docente.

En cumplimiento de la Ley Orgánica de Protección de Datos Personales \cite{lopdp2021} y las normas de entrega de la asignatura, los datos personales de terceros, las grabaciones de audio y los consentimientos firmados no se publican en el repositorio público de GitHub; se conservan de forma privada en una carpeta institucional (\texttt{@uteq.edu.ec}) compartida con el docente.

\subsection{Identificación de la grabación y manifiesto de evidencias}

\begin{tabularx}{\textwidth}{
    |>{\columncolor{azulclaro}\raggedright\arraybackslash}p{4.4cm}
    |>{\raggedright\arraybackslash}X|
}
\hline
\rowcolor{azul}
\textcolor{white}{\textbf{Campo}} &
\textcolor{white}{\textbf{Detalle}} \\
\hline
Código de evidencia &
AUD-01 \\
\hline
Archivo &
\texttt{GA-TL-02\_P01\_audio.m4a} \\
\hline
Formato &
M4A (5,9 MB) \\
\hline
Fecha del ensayo &
20 de noviembre de 2026 \\
\hline
Hora de inicio declarada &
20:00 \\
\hline
Modalidad &
Llamada virtual mediante Zoom \\
\hline
Duración de la grabación &
9 minutos y 24 segundos \\
\hline
Entrevistador &
Kevin German Contreras Chavez \\
\hline
Entrevistado &
Angelo Maximiliano Espinoza Pivaque \\
\hline
Rol representado &
Estudiante que solicita tutorías académicas \\
\hline
Código del participante &
P01 \\
\hline
Carpeta privada institucional &
\url{https://drive.google.com/drive/folders/1tWGhOfK9iarSJlFR9KidZJLUamUjjHuu?usp=sharing} \\
\hline
Acceso del docente &
\texttt{gguerrero@uteq.edu.ec} \\
\hline
Huella SHA-256 &
\texttt{\scriptsize fbf53c7aae8a7d2bd40e73945c0fd38ad046d37abe5b600e1f78b967b067fa58} \\
\hline
\end{tabularx}

\subsection{Alcance del registro}

El archivo de audio se identifica como evidencia
del ensayo. La síntesis de respuestas y los demás
artefactos de análisis corresponden al escenario
representado durante el juego de roles.

Los datos de identificación de la grabación permiten
relacionarla con el informe. Su ubicación privada,
los permisos de acceso, el registro del consentimiento
y el hash SHA-256 quedan consolidados en el manifiesto
de evidencias.

\section{Síntesis de respuestas}

\subsection{Solicitud y coordinación de las tutorías}

El participante indicó que inicia la solicitud de tutoría
mediante una conversación presencial con el docente al
finalizar la clase. En ese contacto explica sus dudas
y consulta la posibilidad de recibir apoyo. Posteriormente,
utiliza WhatsApp para acordar los detalles de la sesión.

La selección del horario depende de la disponibilidad del
docente y de las actividades académicas del estudiante.
Cuando necesita modificarlo, el participante comunica
la situación y espera una nueva confirmación antes de
considerar establecido el cambio.

\subsection{Confirmación e información previa}

El participante considera que una tutoría está confirmada
cuando el docente acepta expresamente el día y la hora.
Una propuesta de disponibilidad no equivale, para él,
a una confirmación.

Cuando se plantean distintos horarios, revisa los mensajes
para identificar el acuerdo definitivo. Si persisten dudas,
consulta nuevamente al docente. Antes de asistir, necesita
conocer el horario, el lugar y los ejercicios o materiales
que debe preparar.

\subsection{Conservación y recuperación de información}

Durante la tutoría, el participante toma apuntes y conserva
los ejercicios realizados. También guarda los archivos
compartidos por el docente para utilizarlos posteriormente.

Para retomar un tema, consulta el cuaderno, los archivos
y los mensajes anteriores. Cuando no encuentra una
recomendación o un documento, necesita solicitar nuevamente
la información al docente.

\subsection{Dificultades y necesidades expresadas}

La principal dificultad señalada es recuperar información
distribuida entre diferentes medios. El participante necesita
consultar sus apuntes y conversaciones para recordar
lo trabajado y las actividades que debe practicar.

Respecto a la duración, explicó que su aprovechamiento
depende de la complejidad del tema y de la oportunidad
de resolver ejercicios. Considera necesario practicar
durante la sesión para comprobar su comprensión.

Finalmente, expresó interés en terminar cada tutoría
con actividades claras y revisar su cumplimiento en
la siguiente sesión, para reconocer las dificultades
que todavía requieren apoyo.

\subsection{Alcance de los resultados}

La actividad corresponde a un juego de roles académico.
El participante confirmó que utilizó respuestas preparadas
previamente con apoyo de inteligencia artificial. Por ello,
el contenido permite analizar el escenario representado,
pero no constituye evidencia independiente de las prácticas
o necesidades de la comunidad estudiantil.

El participante manifestó que comprendió todas las preguntas
y no identificó dificultades de interpretación.

\section{Mapa de empatía}

El mapa organiza en cuatro cuadrantes las expresiones y conductas descritas
por el participante durante la simulación \cite{pohl2025, ireb2026}. Los apartados
sobre pensamientos y sentimientos contienen interpretaciones
del equipo y se distinguen de las afirmaciones expresas.
No se atribuyen conductas observadas fuera del ensayo.

\begin{tabularx}{\textwidth}{
    |>{\columncolor{azulclaro}\raggedright\arraybackslash}X
    |>{\raggedright\arraybackslash}X|
}
\hline
\rowcolor{azul}
\textcolor{white}{\textbf{¿Qué dice?}} &
\textcolor{white}{\textbf{¿Qué piensa?}} \\
\hline
Solicita apoyo al docente después de la clase y continúa la coordinación mediante WhatsApp. \newline
Necesita una confirmación expresa para considerar acordada la fecha y la hora. \newline
Consulta apuntes, archivos y mensajes para recuperar lo trabajado. \newline
Necesita resolver ejercicios para comprobar si comprendió la explicación. \newline
Le resulta útil conocer qué debe practicar y revisar esas actividades en la siguiente sesión. &
Como interpretación de las respuestas, el estudiante representado considera importante distinguir una propuesta de horario de una confirmación definitiva. \newline
También relaciona el aprovechamiento de la tutoría con la posibilidad de practicar y retomar los ejercicios posteriormente. \newline
Estas interpretaciones se apoyan en sus explicaciones y no representan pensamientos observados directamente. \\
\hline
\rowcolor{azul}
\textcolor{white}{\textbf{¿Qué hace?}} &
\textcolor{white}{\textbf{¿Qué siente?}} \\
\hline
Se acerca al docente para comunicar sus dudas. \newline
Coordina horarios y cambios mediante mensajes. \newline
Revisa la conversación para confirmar acuerdos. \newline
Toma apuntes y guarda ejercicios y archivos. \newline
Busca información anterior o vuelve a consultar al docente cuando no logra recuperarla. &
El participante no expresó directamente emociones concretas. \newline
A partir de su necesidad de confirmar horarios y comprobar lo aprendido, el equipo interpreta una búsqueda de seguridad respecto a los acuerdos y a su comprensión de los temas. \newline
No se cuenta con información suficiente para afirmar sentimientos específicos como frustración o ansiedad. \\
\hline
\end{tabularx}

\section{Definición del problema}

\subsection{Usuario}

Estudiante representado en el ensayo que solicita tutorías
académicas para resolver dudas y reforzar su aprendizaje.

\subsection{Necesidad}

Consultar los acuerdos de las tutorías y recuperar los
ejercicios, recomendaciones y actividades pendientes,
para dar continuidad al aprendizaje entre sesiones.

\subsection{Hallazgo}

En el escenario representado, la información necesaria
para continuar el acompañamiento se encuentra distribuida
entre conversaciones, apuntes y archivos. Aunque el estudiante
conserva información de las sesiones, recuperarla requiere
consultar diferentes medios y, en ocasiones, volver
a preguntar al docente.

Además, comprender una explicación durante la tutoría
no basta para comprobar el aprendizaje: el participante
necesita practicar y revisar posteriormente sus dificultades.

\subsection{Enunciado del problema}

El estudiante que recibe tutorías necesita recuperar con
claridad los acuerdos y las actividades de cada sesión,
porque la distribución de la información entre distintos
medios dificulta retomar lo trabajado y reconocer qué
aspectos requieren apoyo adicional.

Este enunciado se deriva de la simulación y deberá
contrastarse con participantes reales antes de utilizarlo
como diagnóstico del proyecto.

\subsection{Pregunta: ¿Cómo podríamos...?}

¿Cómo podríamos facilitar que los estudiantes recuperen
los acuerdos y las actividades de sus tutorías para
dar continuidad a su aprendizaje e identificar las
dificultades que requieren apoyo adicional?

\section{Revisión y ajustes del guion}

\subsection{Claridad durante el ensayo}

El participante manifestó que comprendió todas las preguntas
y no señaló dificultades de interpretación. Por tanto,
no se identificaron ajustes motivados por confusión
declarada durante el ensayo.

Esta valoración corresponde a una aplicación con respuestas
preparadas previamente, por lo que su alcance para evaluar
la comprensión espontánea del guion es limitado.

\subsection{Mejoras identificadas en la revisión posterior}

La revisión del guion permite proponer mejoras de precisión
y cobertura. Las siguientes observaciones corresponden
al análisis del instrumento y no a dificultades comunicadas
por el participante.

\begin{description}
    \item[Preguntas 1 y 2: posible solapamiento.]
    La explicación de cómo se solicita una tutoría puede
    incluir el medio de contacto, que vuelve a consultarse
    en la segunda pregunta. Para obtener información
    complementaria, se propone reformular la pregunta 2:

    \emph{¿Qué dificultades encuentras al comunicarte
    con el docente para solicitar una tutoría?}

    \item[Pregunta 4: precisión de la situación.]
    La formulación original no distingue quién propone
    el cambio. Para centrar la respuesta en la experiencia
    del estudiante, se propone:

    \emph{Cuando necesitas cambiar el horario de una
    tutoría acordada, ¿qué pasos sigues?}

    \item[Pregunta 12: cobertura del enfoque actualizado.]
    El guion se concentra en la coordinación de tutorías
    y aborda de forma limitada la identificación temprana
    de dificultades académicas. Se propone sustituir
    la pregunta sobre duración por:

    \emph{¿Qué señales te permiten reconocer que necesitas
    apoyo académico antes de una evaluación?}
\end{description}

\subsection{Aplicación de los ajustes}

Las reformulaciones son propuestas para una aplicación
posterior. No se presentan como preguntas realizadas
durante el ensayo ni se les atribuyen respuestas.

La versión aplicada se conserva en el guion para mantener
la correspondencia con la grabación. Una nueva aplicación
permitiría evaluar las preguntas revisadas y explorar
aspectos que este ensayo no cubrió.

\section{Declaración de uso de inteligencia artificial}

En cumplimiento de las normas de integridad académica de la asignatura y el código ético de la ACM \cite{acm2018}, el equipo declara el apoyo recibido de herramientas de inteligencia artificial para este ensayo:

\begin{tabularx}{\textwidth}{
    |>{\columncolor{azulclaro}\raggedright\arraybackslash}p{4.1cm}
    |>{\raggedright\arraybackslash}X|
}
\hline
\rowcolor{azul}
\textcolor{white}{\textbf{Campo}} &
\textcolor{white}{\textbf{Detalle}} \\
\hline
Herramienta y versión &
ChatGPT y asistente de redacción/código (versión utilizada durante la preparación del ensayo, noviembre de 2026). \\
\hline
Propósito &
Apoyo en la redacción de respuestas preparadas para el juego de roles \texttt{P01}, estructuración y depuración de sintaxis LaTeX, y organización del borrador del informe. \\
\hline
Verificación realizada &
El equipo revisó manualmente cada sección frente a la guía \texttt{GA-TL-02-FORM}, cotejó las afirmaciones con la grabación de audio \texttt{GA-TL-02\_P01\_audio.m4a} y compiló el documento en LaTeX para verificar su correcta presentación. \\
\hline
\end{tabularx}

\begin{thebibliography}{9}

\bibitem{pohl2025}
K. Pohl, \emph{Requirements Engineering: Fundamentals, Principles, and Techniques}, 2nd ed. Berlin, Heidelberg, Germany: Springer, 2025, ISBN: 978-3-662-69204-2.

\bibitem{ireb2026}
S. B\"uhne and M. Glinz, \emph{IREB Certified Professional for Requirements Engineering -- Foundation Level Syllabus, Version 3.3.0}. Karlsruhe, Germany: International Requirements Engineering Board (IREB) e.V., Apr. 2026.

\bibitem{lopdp2021}
Asamblea Nacional del Ecuador, \emph{Ley Org\'anica de Protecci\'on de Datos Personales}, Registro Oficial, Quinto Suplemento No. 459, 26 de mayo de 2021.

\bibitem{acm2018}
Association for Computing Machinery (ACM), \emph{ACM Code of Ethics and Professional Conduct}. 2018. [Online]. Available: \url{https://www.acm.org/code-of-ethics}

\bibitem{chang2025}
Y.-H. Chang, F.-C. Chen, and C.-I. Lee, ``Developing an Early Warning System with Personalized Interventions to Enhance Academic Outcomes for At-Risk Students in Taiwanese Higher Education,'' \emph{Education Sciences}, vol. 15, no. 10, p. 1321, Oct. 2025, doi: \url{https://doi.org/10.3390/educsci15101321}.

\bibitem{martins2019}
H. Ferreira Martins, A. Carvalho de Oliveira Junior, E. Dias Canedo, R. A. Dias Kosloski, R. \'Avila Vald\^es, and E. Costa Oliveira, ``Design Thinking: Challenges for Software Requirements Elicitation,'' \emph{Information}, vol. 10, no. 12, p. 371, Nov. 2019, doi: \url{https://doi.org/10.3390/info10120371}.

\end{thebibliography}

\end{document}
